% =====================================================================
% Apartado 7: Comparación y postura crítica (1 min) -- Responsable: Elian Mendoza
% Codex asistió en las precisiones para revisión de Elian.
% =====================================================================
\section{Comparación y postura crítica}

% ---------------------------------------------------------------------
\begin{frame}{Comparación con SoundCloud: diseño y consistencia}
  \scriptsize
  \renewcommand{\arraystretch}{1.3}
  \begin{tabular}{p{0.16\textwidth} p{0.37\textwidth} p{0.37\textwidth}}
    \toprule
    \textbf{Criterio} & \textbf{Spotify} & \textbf{SoundCloud} \\
    \midrule
    Arquitectura &
      Unos 1.200 servicios y flujos de datos migrados a Google Cloud
      \cite{googleSpotify} &
      Microservicios con capa BFF; sustituyó un monolito en Ruby on
      Rails \cite{obrien2022strangler} \\
    Almacenamiento &
      Cassandra para personalización \cite{mishra2015personalization} &
      Cassandra \cite{pythianSoundcloud}; Roshi, caso histórico sobre
      Redis \cite{bourgon2014roshi,soundcloudRoshi} \\
    Consistencia &
      Configurable por operación: ONE o QUORUM \cite{cody2015cassandra} &
      Roshi: CRDT de tipo LWW, quórum de escritura y reparación en
      algunas estrategias de lectura \cite{almeida2024crdt,soundcloudRoshi} \\
    CAP &
      Padlock: Cassandra de alta disponibilidad con lecturas LOCAL\_ONE
      por región \cite{leenders2018padlock} &
      Roshi: tolerante a particiones, altamente disponible y
      eventualmente consistente \cite{soundcloudRoshi} \\
    PACELC &
      La ruta local de Padlock y la personalización favorecen baja latencia
      \cite{mishra2015personalization,leenders2018padlock} &
      Roshi permite converger sin consenso para ordenar las actualizaciones.
      Las escrituras esperan un quórum \cite{soundcloudRoshi} \\
    \bottomrule
  \end{tabular}
\end{frame}

% ---------------------------------------------------------------------
\begin{frame}{Comparación con SoundCloud: resultados y compromisos}
  \scriptsize
  \renewcommand{\arraystretch}{1.3}
  \begin{tabular}{p{0.16\textwidth} p{0.37\textwidth} p{0.37\textwidth}}
    \toprule
    \textbf{Criterio} & \textbf{Spotify} & \textbf{SoundCloud} \\
    \midrule
    Disponibilidad &
      Padlock: tres regiones con tres zonas; objetivo de 99,95~\%
      \cite{leenders2018padlock} &
      Cassandra en varios centros de datos \cite{pythianSoundcloud};
      Roshi, normalmente con tres clústeres \cite{soundcloudRoshi} \\
    Latencia medida &
      Menos de 5~ms en el percentil 95 con seis nodos
      \cite{cody2015cassandra} &
      Lecturas en el percentil 99: de 250~ms a 40~ms
      \cite{pythianSoundcloud} \\
    Fortaleza &
      Nivel de consistencia ajustable por servicio
      \cite{cody2015cassandra} &
      Roshi: escrituras rápidas y almacenamiento reducido
      \cite{bourgon2014roshi} \\
    Limitación &
      Con tres nodos no alcanzó el objetivo de latencia
      \cite{cody2015cassandra} &
      Roshi no garantiza durabilidad completa y depende de marcas de
      tiempo correctas \cite{soundcloudRoshi} \\
    Compromiso &
      Mayor consistencia mediante niveles de lectura más exigentes, a
      cambio de mayor latencia
      \cite{cody2015cassandra,leenders2018padlock} &
      Escrituras rápidas y almacenamiento reducido, a cambio de lecturas
      más complejas y consistencia eventual \cite{bourgon2014roshi} \\
    \bottomrule
  \end{tabular}
\end{frame}

% ---------------------------------------------------------------------
\begin{frame}{Postura crítica}
  \small
  Los casos documentados de personalización y Roshi responden a
  necesidades distintas.

  \bigskip
  \renewcommand{\arraystretch}{1.4}
  \begin{tabular}{p{0.18\textwidth} p{0.34\textwidth} p{0.34\textwidth}}
    \toprule
    & \textbf{Spotify} & \textbf{SoundCloud (Roshi)} \\
    \midrule
    Mecanismo & Ajusta cuántas réplicas deben responder & Converge
      con un CRDT al recibir las mismas actualizaciones \\
    Conveniente\newline cuando & Distintos datos requieren distintos
      niveles de consistencia & La prioridad es no bloquear las
      operaciones \\
    Riesgo & Mayor latencia al exigir más réplicas & Sin durabilidad
      perfecta: perder réplicas antes de repararlas puede perder datos \\
    \bottomrule
  \end{tabular}

  \bigskip
  La elección depende de qué retraso y pérdida de datos tolera cada flujo.
  Las cifras de latencia proceden de cargas y percentiles distintos.
\end{frame}
